\problema{Clone Graph}{133}{src/01-grafos/133-clone-graph.cpp}{M}

Nos dan una referencia a un nodo de un grafo no dirigido y \emph{conexo}.
Cada nodo tiene un valor entero (\texttt{val}) y una lista de vecinos
(\texttt{neighbors}). Queremos devolver una \textbf{copia profunda} de todo
el grafo: cada nodo de la copia debe ser un objeto nuevo, con el mismo valor
y las mismas conexiones que el original.

\paragraph{La idea, con un ejemplo de la vida real.} Imagina que tienes un
grupo de amigos dibujado con líneas que unen a quienes se conocen, y quieres
hacer una \emph{copia nueva} del grupo entero: personas nuevas, de carne y
hueso, pero con los mismos nombres y las mismas amistades. A eso se le llama
\textbf{copia profunda}: no vale con señalar a los amigos originales, hay que
crear a cada uno de nuevo.

Recorremos el grupo con un \emph{DFS} (ir metiéndonos por un amigo, luego por
un amigo de ese amigo, y así hasta el fondo antes de regresar). El único
problema es que las amistades forman \emph{ciclos}: A conoce a B, y B conoce a
A. Si no tenemos cuidado, clonaríamos a A, luego a B, luego a A otra vez,
luego a B otra vez\ldots para siempre. La solución es llevar una
\emph{libreta} que anota, para cada persona original, cuál es su copia ya
hecha. Antes de recorrer los amigos de alguien, apuntamos su copia en la
libreta; así, si volvemos a toparnos con él, simplemente usamos la copia que
ya está anotada en vez de crearla de nuevo.

\begin{center}
\ttfamily
\begin{tabular}{ccc}
1 & --- & 2\\
$|$ & & $|$\\
4 & --- & 3\\
\end{tabular}
\end{center}

\noindent Lista de adyacencia de este ejemplo: $1:[2,4]$, $2:[1,3]$,
$3:[2,4]$, $4:[1,3]$. Al clonar desde el nodo $1$, el DFS visita $2$, que a
su vez intenta visitar $1$ otra vez: como $1$ ya está en el mapa, se
reutiliza su copia en lugar de recrearla.

\paragraph{Cómo lo resuelve el código, paso a paso.} Usamos un mapa
\texttt{unordered\_map<Node*, Node*>} como esa libreta: dado un nodo
original, te dice cuál es su copia. La función recursiva, al llegar a un
nodo, hace:
\begin{enumerate}
  \item Si el nodo ya está en la libreta, devolvemos directamente su copia
        (ya lo clonamos antes). Esto es lo que rompe los ciclos y evita la
        recursión infinita.
  \item Si no está, creamos su copia (mismo valor) y la anotamos en la
        libreta \emph{de inmediato}, antes de tocar a sus vecinos.
  \item Recorremos cada vecino del original, lo clonamos con la misma
        función y vamos agregando esas copias a la lista de vecinos de
        nuestra copia.
  \item Devolvemos la copia ya con todos sus vecinos conectados.
\end{enumerate}

\lstinputlisting{src/01-grafos/133-clone-graph.cpp}

\complejidad{$O(V+E)$}{$O(V)$}

\paragraph{¿Qué significa esa complejidad?} $V$ es el número de nodos
(personas) y $E$ el de aristas (amistades). $O(V+E)$ significa que tocamos
cada persona una vez y recorremos cada amistad una vez: no se puede copiar el
grupo sin mirarlo entero al menos una vez. El espacio $O(V)$ es por la
libreta (una entrada por persona) más la pila de la recursión.

\paragraph{Consejos para la entrevista (Amazon).} El error típico es anotar
la copia en la libreta \emph{después} de recorrer a los vecinos: eso provoca
recursión infinita en cuanto hay un ciclo (y casi siempre lo hay, porque el
grafo es no dirigido). Menciona también la versión con BFS y una cola (la
misma idea de la libreta, pero sin recursión, útil si te preocupa que la pila
se llene en grafos enormes). Casos raros: nodo nulo (grafo vacío) devuelve
\texttt{nullptr}, un solo nodo sin vecinos, y comprueba que las copias sean
objetos de verdad nuevos (direcciones de memoria distintas a las del
original), no un simple ``apodo'' que apunta al mismo lugar.
